% TOP_PROBLEM: {"problemId": "problem.amenability-of-thompson-s-group-f", "canonicalTitle": "Amenability of Thompson's Group F", "releaseRank": 304, "primaryDomain": "algebra_representation_category", "primaryDomainLabel": "Algebra, representation & category theory", "displayStatus": "open", "releaseStatus": "open"}
% !TeX program = lualatex
% Definition and sources only; this file is not an LLM solution attempt.
\documentclass[11pt]{article}
\usepackage[margin=25mm]{geometry}
\usepackage{fontspec}
\setmainfont{Latin Modern Roman}
\usepackage{amsmath,amssymb,mathtools}
\usepackage{unicode-math}
\setmathfont{Latin Modern Math}
\usepackage{xurl,hyperref}
\hypersetup{colorlinks=true,urlcolor=blue}
\setcounter{secnumdepth}{0}
\setlength{\parindent}{0pt}
\setlength{\parskip}{5pt}
\setlength{\emergencystretch}{3em}
\begin{document}
\section*{\texorpdfstring{Amenability of Thompson's Group F}{Amenability of Thompson's Group F}}\label{304}
\subsection{Definitions and mathematical statement}
Thompson's group $F$ consists of increasing piecewise-linear homeomorphisms of $[0,1]$ with finitely many breakpoints, all dyadic rational, and slopes in $\{2^k:k\in{\mathbb{Z}}\}$; its operation is composition. A discrete group is amenable if it has a left-invariant mean: a positive linear functional $m:\ell^\infty(F)\to{\mathbb{R}}$ with $m(1)=1$ and $m(f\circ L_g)=m(f)$ for every $g\in F$. Equivalently, its finite subsets satisfy the F\o lner criterion
\[
 \forall E\subset F\text{ finite}\ \forall{\varepsilon}>0\ \exists A\subset F,
 \quad0<|A|<\infty,\quad |gA\mathbin{\triangle} A|<{\varepsilon}|A|\ (g\in E).
\]
Determine whether $F$ is amenable. Absence of a nonabelian free subgroup alone is not the definition and would not settle this question.
\par\smallskip\textit{Formulation and concept references:} \href{https://arxiv.org/abs/2603.28228}{[S1]}.

\subsection{Short English statement}
Does Thompson's group admit an averaging operation invariant under all translations, or equivalently finite sets with arbitrarily small translation boundaries?
\subsection{Sources}
\begin{itemize}
\item[S1]Anna Cascioli, Martin Gilabert Vio, Eduardo Silva, Stationary boundaries on the space of amenable subgroups and C*-simplicity, arXiv:2603.28228 (2026).
\url{https://arxiv.org/abs/2603.28228}.
\textit{Formulation source.}
\end{itemize}
\end{document}
